\documentstyle[a4r,fancybox,amstex,graphicx,semrot,epsf,psfig,semcolor,slidesec]{seminar}

\slidesmag{3}                           % magstep definition for slides
\articlemag{1}                          % magstep definition for article st
%\epsfslidesize                          % auto scaling of postscript files
%\portraitonly                          % To print only portrait slides
%\landscapeonly                         % To print only landscape slides
\rotateheaderstrue                     % Try this out if using rotation ma
\input amssym.def
\input amssym.tex
%\input amsmath.sty

\setlength\slidewidth{10in}\setlength\slideheight{6.9in}

%\setlength\slidewidth{6.9in}\setlength\slideheight{10in}

\renewcommand{\slidetopmargin}{5.7in}
\renewcommand{\slidebottommargin}{0in}
\renewcommand{\slideleftmargin}{4.5in}
\renewcommand{\sliderightmargin}{1in}

\newcommand{\minimize}{\mathop{minimize}\limits_{\bf X \epsilon{\cal R}^n}}

\def\printlandscape{\special{landscape}}    % Works with dvips

\newcommand{\heading}[1]{\begin{center}\large\bf \shadowbox{#1}\end{center}
                         \vspace{1ex minus 1ex}}

\newcommand{\bigheading}[1]
           {\slideheading{#1}
\begin{center}
\shadowbox{\begin{minipage}{8.0truein}
\begin{center}\large\bf
#1\end{center}\end{minipage}}
\end{center}
           \vspace{1ex minus 1ex}}



\newpagestyle{icim}
     {}
     {\black \hspace{.99in}\hspace{\fill} \theslide}%

\pagestyle{icim}
\slideframe{oval}
\def\Reel{\mbox{I\hspace{-.20em}R}}


\newcommand{\dessin}[4]{
\begin{figure}[H]
\centering
\HideDisplacementBoxes
\BoxedEPSF{#1 scaled #2}
\caption{#3}
\label{#4}
\end{figure}}

\def\text{\hbox}
\def\qed{\unskip\kern 10 pt\penalty 500 
\raise -2pt\hbox{\vrule\vbox to9pt{\hrule width 7pt
\vfill\hrule}\vrule}}
\def\itt#1{{\it (#1)\/}}
\def\dd#1#2{\frac{d#1}{d#2}}
\def\ddd#1#2{\frac{d^2{#1}}{d{#2}^2}}
\def\ddq#1#2{\frac{d^4{#1}}{d{#2}^4}}
\def\ddj#1#2{\frac{d^j{#1}}{d{#2}^j}}
\def\dpa#1#2{\frac{\partial #1}{\partial #2}}
\def\ddpa#1#2{\frac{\partial^2 {#1}}{\partial {#2}^2}}
\def\ddpab#1#2#3{\frac{\partial^2 #1}{\partial #2 \partial #3}}
\def\addots{\mathinner{\mkern1mu\raise1pt\hbox{.}\mkern2mu\raise4pt
\hbox{.}\mkern2mu\raise7pt\vbox{\kern7pt\hbox{.}}\mkern1mu}}
\def\Ki{{\cal X}}
\def\epsi{\varepsilon}
\def\grad{\overrightarrow{\text {grad}}}
\def\dive{\text {div}\,}
\def\tg{\text {tg }}
\def\cotg{\text {cotg }}
\def\rot{\text {rot }}
\def\sign{\text {sign }}
\def\curl{\text {curl }}
\def\epi{\text {\'epi }}
\def\imp{\Rightarrow}
\def\equi{\Leftrightarrow}
\def\Ker{\text{Ker }}
\def\Ima{\text{Im }}
\def\supess{\mathop{\text{sup ess}\,}}
\def\Log{\text{Log }}
\def\Supp{\text{Supp }}
\def\supp{\text{supp }}
\def\map\essstow{\rightharpoonup}
\def\tend{\mathop{\longrightarrow}}
\def\notend{\mathop{\not{\!\!\!\!\longrightarrow}}}
\def\union{\mathop{\cup}}
\def\inter{\mathop{\cap}}
\def\Rl{\text {I}\!\text {R}}
\def\Rr{\mathbb R\,}
\def\Cc{\mathbb C\,}
\def\Pp{\mathbb P\,}
\def\RRP{\mathbb R}
\def\RRB{\overline{\mathbb R}\,}
\def\Ne{\text {I}\!\text {N}}
\def\NN{\mathbb N\,}
\def\NNP{\mathbb N}
\def\Zz{\mathbb Z\,}
\def\ZZP{\mathbb Z}
\def\CC{\mathbb C\,}
\def\CCP{\mathbb C}
\def\vc#1{\overrightarrow{#1}}
\def\ti#1{\tilde{#1}}
\def\nab{\nabla}
\def\lap{\Delta}
\def\Lp#1{L^{#1}(\Omega)}
\def\LP#1#2{L^{#1}(0,{T};\-#2)}
\def\C#1#2{{\cal C}^{#1}([0,\text {T}];\-#2)}
\def\Co#1#2{{ C}^{#1}([0,{T}];\-#2)}
\def\LPE#1#2{ L^{#1}(\eta, {T};\-#2)}
\def\LPT#1#2{ L^{#1}(\eta,\infty;\-#2)}
\def\LR#1#2{ L^{#1}(0,\infty;\-#2)}
\def\lp#1{L^{#1}(\Omega)}
\def\h#1{H^{#1}(\Omega)}
\def\cun#1{C^{#1}(\overline{\Omega})}
\def\hc#1{ H^{#1}_0(\Omega)}
\def\w#1#2{W^{#1,#2}(\Omega)}
\def\W#1#2#3{W^{#1,#2}_{#3}((0, T)\times \Omega)}
\def\det{\left|\matrix} \def\enddet{\endmatrix\right|}
\def\dete{\text{d\'et }}
\def\dsum{\displaystyle\sum}

\begin{document}
 
 
 
 
 % transp no 1
\begin{slide}
\pagestyle{empty}
\vspace{1cm}
{\LARGE \bf \begin{center}  
Numerical properties of currents associated to
transcendental entire curves
 \end{center}}
\vspace{1cm}
{\large \bf \begin{center}  
{Mihai P\u aun}
 \end{center}}
\vspace{2cm}
\begin{center}

$\strut$ \llap Universit\'e Louis Pasteur, Strasbourg
\end{center}
\end{slide}

% transp no 2

\begin{slide}

\par\noindent{\large \bf Motivation}
\par\vspace{\bigskipamount}\noindent
\par\noindent We recall the following theorem
\vskip 0.3cm
\noindent{\bf Theorem}{ (McQuillan, 98)} 

{\sl Let $S$ be a projective surface,
endowed with a holomorphic foliation ${\cal F}$. 

Let $\varphi:{\Cc} \mapsto S$ be a holomorphic entire curve, tangent to ${\cal F}$.

Then $\varphi$ is algebraically degenerate.}


\vskip 0.5cm

$\bullet$ $\varphi \mapsto T[\varphi]$

\vskip 3pt

$\bullet$ $(T[\varphi], N_{\cal F}^*)$ is bounded by the value of the 
logarithmic singularities of 
 $T[\varphi]$ at ${\cal F}_s$

\vskip 3pt

$\bullet$ Aim: present some of the numerical properties of $T[\varphi]$ in higher 
dimensional setting, in the \vskip 0cm
\hskip 0.3cm case when $T[\varphi]$ acquires 
singularities at points or analytic subsets
of the ambient manifold.
\end{slide}




\begin{slide}
\noindent {\large\bf Nef cohomology classes}

\vskip 10pt
\par\noindent 
{\bf Definition} {\sl Let $X$ be a compact complex manifold, endowed with a K\"ahler metric
$\omega$  

\hskip 1.8cm A cohomology class $\{\alpha\}\in H^{1, 1}(X, {\Rr})$ is nef if } 
\vskip 5pt
\centerline {\sl   
$\exists \hspace{0.1cm} (f_\varepsilon)_{\varepsilon> 0}\subset {\cal C}^{\infty}(X)$ \hskip 0.3cm
s. t. \hskip 0.3cm
$\alpha+ i\partial\overline\partial f_\varepsilon\geq -\varepsilon \omega$ \hskip 0.3cm on 
\hskip 0.3cm$X$}
\vskip 5pt 
\hskip 1.8cm The {\sl numerical dimension} of a nef class ${\alpha}$ is defined by 
$$\nu(\alpha):= \max \{ k / \int_X\alpha^k\wedge\omega^{n- k}> 0\}$$

{\bf Basic properties }

\par\vspace{\bigskipamount}\noindent
$\bullet$ If $\{\alpha\}= c_1(L)$, and if $X$ is projective (or even Moishezon) then 
\vskip 3pt 
\centerline{ $\{\alpha\}$ is nef
  $\iff$ $L$ is nef in the usual sense}
\vskip 3pt
\par\vspace{\bigskipamount}\noindent
$\bullet$ $\{\alpha\}$ is nef $\iff$ $\int_Y\alpha^k\wedge \omega^{p-k}\geq 0$ \hskip 0.3cm 
$\forall$ compact analytic sets $Y$ and all $k$
\vskip 3pt
\par\vspace{\bigskipamount}\noindent
$\bullet$ In general, a nef class does not contain a smooth and semi--positive representative,

\hskip 0.3cm it only has a weak positivity property

\end{slide}


\begin{slide}
\noindent {\large\bf Closed positive currents}

\vspace{0.5cm}
\par\noindent 
$\bullet$ A {\sl current} $T$ is a differential form with distribution coefficients; 
\vskip 3pt
$\rightarrow$ in the complex situation, we are interested in currents of pure bidegree $(p,p)$
$$
T=i^{p^2}\sum_{|I|=p,|J|=p} T_{I,J}\,dz_I\wedge d\overline z_J\qquad
\hskip 1cm
\left|
\begin{array}{ll}
T_{I,J}\mbox{ distributions on }X,\\ \\
\mbox{here }dz_I:=dz_{i_1}\wedge\ldots\wedge dz_{i_p} 
\mbox{ as usual.}\end{array}\right. 
$$
\vskip 5pt
\par\noindent
$\bullet$ {$T$ is closed if $\langle T, du\rangle= 0$}
\vskip 5pt
 \par\noindent
$\bullet$ We say that $T$ is positive if $ \hskip 0.5cm
\left|
\begin{array}{ll}
T= \overline T \\
\mbox{\hskip 0.5cm}\& \\
\sum \lambda_I\overline\lambda_JT_{I,J}\mbox{ is a positive distribution }
\forall\lambda_I\in {\CC},\ |I|=p
\end{array}\right. 
$
\vskip 5pt
 \par\noindent
$\bullet$ The Lelong number of $T$ at the point $x_0\in X$ is defined as 
$$\nu (T, x_0):= \lim_{r\mapsto 0}{{1}\over {r^{2(n- p)}}}\int_{\vert w\vert < r}T\wedge 
(i\partial\overline\partial \vert w\vert ^2)^{n-p}$$

\end{slide}

\begin{slide}
\noindent {\large \bf Examples}
\par\vspace{\bigskipamount}\noindent
\par\noindent
$\bullet$ Let $Y\subset X$ be a analytic set of codimension $p$ ;
its associated current $T=[Y]$ is defined by
$$
\langle[Y], u\rangle=\int_Y u,\qquad u\in {\cal D}_{n-p,n-p}(X)
$$
\hskip 0.5cm $\rightarrow$ $[Y]$ is a closed positive current of bidegree $(p,p)$.

\par\vspace{\bigskipamount}\noindent
$\bullet$ Let $\Omega\subset {\CC}^n$ be an open set, and let $(f_j)\subset {\cal O}(\Omega)$
 be a set of holomorphic functions on $\Omega$ ; then 
\vskip 3pt 
\hskip 1cm $T= i\partial\overline \partial\log(\vert f_1\vert ^2+ ...+\vert f_k\vert ^2)$ \hskip 1cm
is a closed positive current of bidegree $(1,1)$. 

\par\vspace{\bigskipamount}\noindent
$\bullet$ Let $L\mapsto X$ be a line bundle, $h$ a hermitian metric on $L$ ;  
$\sigma\in H^0(X, L)$ a non-zero section. 
\vskip 3pt
\hskip 1cm
$T= \Theta_h(L)+ i\partial\overline \partial\log(\vert \sigma\vert_h^2)$
\hskip 1cm
 is a closed positive current
in $c_1(L)$ 
\vskip 1pt 
\hskip 6cm $\rightarrow$ and it equals $[\sigma= 0]$

\par\vspace{\bigskipamount}\noindent
$\bullet$ 
Let $\{\beta \}\in H^{p,p}(X, {\Rr})$, and take $(\gamma_j)\in \{\beta \}$ such that
$\gamma_j\geq 0$, for all $j$. Then :
 \vskip 3pt
\centerline{ $\exists\ \Theta\in \{\beta \}$, such that 
$\Theta= \lim \gamma_{j_k}$, in weak sense on $X$} 
%i.e. 
%$\langle\Theta, u\rangle= \lim\langle \gamma_{j_k}, u\rangle $.

\par\vspace{\bigskipamount}\noindent
$\bullet$ Let $\{\alpha \}\in H^{1,1}(X, {\Rr})$ be a nef class ;
 then there exist
a closed positive current $T\in \{\alpha \}$ 
\end{slide}

\begin{slide}
\noindent {\Large\bf Concentration of mass for nef classes}

\par\vspace{\bigskipamount}\noindent
{\bf Theorem 1} {\sl Let $(X, \omega)$ be a compact K\"ahler manifold, and let 
$\{\alpha \}\in H^{1,1}(X, {\Rr})$ be a nef 

\hskip 1.9cm cohomology class.
If $Y$ is a codimension $p$ analytic subset of $X$, then}
\par\vspace{\bigskipamount}\noindent
\centerline{\sl $\forall \varepsilon > 0$, $\exists$
$T_\varepsilon\in \{\alpha+ \varepsilon\omega\}$ \hskip 0.2cm s. t.  \hskip 0.2cm
$T_\varepsilon\geq 0$ \hskip 0.2cm and \hskip 0.2cm 
$\nu_Y(T_\varepsilon)\geq C \varepsilon^{{{n- \nu(\alpha)}\over {p}}}$}
\par\vspace{\bigskipamount}\noindent
\hskip 1.9cm {\sl for some positive constant $C$ independant of $\varepsilon$.}

\par\vspace{\bigskipamount}\noindent
{\bf Remark.} If $(X, A)$ is a polarised algebraic manifold, 
\hskip 0.2cm $\{\alpha\}= c_1(L)$ \hskip 0.2cm and \hskip 0.2cm $Y= x_0$

$\bullet$ $\forall m$, the bundle $mL+ A$ is ample, so by Riemann-Roch 
$$h^0\big(X, k(mL+ A)\big)\geq
Ck^nm^{\nu(L)}$$


$\bullet$ $\exists$
$\sigma\in H^0\big(X, k(mL+ A)\big)$ \hskip 0.2cm such that \hskip 0.2cm 
$ord_{x_0}(\sigma)\geq km^{\nu(L)/n}$\hskip 0.2cm
(by linear algebra)


$\bullet$ $\displaystyle T_k:= {{1}\over {mk}}[Z_\sigma]$  \hskip 0.2cm 
where $[Z_\sigma]$ is the zero set of 
the section $\sigma$.

\end{slide}


\begin{slide}

\noindent{\large\bf Construction of $T_\varepsilon$ ; the case $\alpha= \omega$}
\par\vspace{\bigskipamount}\noindent
{\bf  Local model for $T_\varepsilon$ }:
 consider $U\subset X$ open set, such that 
$Y\cap U= (g_1= 0, ..., g_p= 0)$. 

\par\noindent The sequence of functions $\psi_{U,\delta}:= 
\log (\displaystyle\sum_{j=1}^p\vert g_j\vert ^2+ \delta^2)$ has the properties :
\vskip 5pt
\par\hspace{0.5cm} $\bullet$  $\psi_{U,\delta}$ is smooth, and 
$i\partial\overline \partial\psi_{U,\varepsilon}\geq 0$.
\vskip 5pt\par\hspace{0.5cm} $\bullet$  the limit 
$i\partial\overline \partial\psi_{U,\delta}
\mapsto T_U$ has logarithmic poles on $Y\cap U$.

\par\vspace{\bigskipamount}\noindent
We {\bf globalize the local construction} by a partition of unit ; define

$$\psi_{\delta}:= \log \bigl(\sum_{j=1}^N\theta_j\exp(\psi_{U_j,\delta})\bigr)$$

\par\hspace{0.5cm} $\bullet$ we have $(\psi_{\delta})\subset {\cal C}^{\infty}(X)$, and 
$\displaystyle \psi:= \lim_{\delta \mapsto 0}\psi_{\delta}$ has logarithmic poles on $Y$.

\vskip 5pt
\par\hspace{0.5cm} $\bullet \hspace{0.1cm} \exists \hspace{0.1cm} c> 0$ such that 
$i\partial\overline \partial\psi_{\delta}\geq -c\omega$, $\forall \delta> 0$
\vskip 5pt
\par\hspace{0.5cm} $\bullet$ the sequence $T_\delta:= \omega+ 
\frac {i}{2c}\partial\overline \partial\psi_\delta\in \{\omega\}$ converge weakly
to a current $T$ 

\par\hspace{0.7cm} and for all $\varepsilon> 0$, we take $T_\varepsilon:= T$.

\end{slide}



\begin{slide}
\noindent{\large\bf Construction of $T_\varepsilon$ ; the general case}

 %We recall the following deep result of S.-T. Yau.
\par\vspace{\bigskipamount}\noindent
{\bf Theorem} {(Yau).} {\sl Let $(X, \omega_0)$ be a compact K\"ahler manifold, and consider
$f$ a smooth volume form on $X$, such that $\displaystyle\int_X\omega_0^n= \int_Xf$. 
Then there exist a
smooth function $\varphi$ such that 
$$\omega_0+ i\partial\overline \partial\varphi> 0\hbox{ and }
(\omega_0+ i\partial\overline \partial\varphi)^n= f.$$}
For each $\varepsilon> 0$, the class $\{\alpha+ \varepsilon\omega\}$ is K\"ahler ; 
\vskip 1pt
$\rightarrow$ take 
the smooth volume form given by the previous construction
$\displaystyle f_\delta:= 
\bigl(\omega+ \frac {i}{2c}\partial\overline \partial\psi_{\delta}\bigr)^n$

\vskip 3pt
Yau's theorem give $(\varphi_{\varepsilon, \delta})\subset{\cal C}^\infty(X)$ such that:
$\left\lbrace
\begin{array}{ll} 
\alpha+  \varepsilon\omega+ 
i\partial\overline \partial\varphi_{\varepsilon, \delta}> 0\mbox{ on }X.\\ \\
\displaystyle (\alpha+  \varepsilon\omega+ 
i\partial\overline \partial\varphi_{\varepsilon, \delta})^n= 
C_\varepsilon f_\delta\end{array}\right.$\vskip 1pt
\hskip 0.5cm where $C_\varepsilon$ is a normalization constant, given by
$\displaystyle C_\varepsilon= \frac{\int_X(\alpha+  \varepsilon\omega)^n}{
\int_X\omega^n}
\geq C\varepsilon^{n- \nu(\alpha)}$, as $\varepsilon\mapsto  0$
\vskip 3pt
Denote by $T_{\varepsilon, \delta}:= \alpha+  \varepsilon\omega+ 
i\partial\overline \partial\varphi_{\varepsilon, \delta}$ ; 
\vskip 1pt
$\rightarrow$ take $T_\varepsilon$ to be 
any weak limit of the sequence $(T_{\varepsilon, \delta})$ 
\vskip 1pt
$\rightarrow$ one can show that 
it has the required properties
\end{slide}

\begin{slide}
\par\vspace{\bigskipamount}\noindent
{\large\bf Currents associated to entire curves}

\par\vspace{\bigskipamount}\noindent
Let $(X, \omega)$ be a compact K\"ahler manifold, and let $\varphi:{\Cc}\mapsto X$
be a non-degenerate entire curve. 
The {\sl characteristic function} of the curve $\varphi$ with respect to
$\omega$ is given by
$$
\displaystyle T(\varphi, r):=  
\int_0^r{{dt}\over {t}}\int_{\Delta(t)}\varphi^*\omega.$$

Classically, Nevanlinna theory gives an equivalent for $T(\varphi, r)$
as follows

$$T(\varphi, r)= \sum _j\nu_j\log{{r}\over {\vert z_j\vert }}+ 
{{1}\over {2\pi}}\int_0^{2\pi}\log{{1}\over {\vert \vert \sigma \circ \varphi (re^{i\theta})
\vert \vert _h}}
d\theta+ {\cal O}(1)$$

where : $\sigma\in H^0(X, A)$\hskip 0.3cm  and \hskip 0.3cm $\Theta_h(A)=\omega$

\hskip 1.3cm $\displaystyle \sum \nu_j\delta_{z_j}$ \hskip 0.2cm is the zero divisor of $\sigma \circ \varphi$

$\bullet$ we define the $(n- 1, n-1)$ current
$\displaystyle (T_r[\varphi], \alpha):= {{\int_1^r{{dt}\over {t}}\int_{\Delta(t)}\varphi^*\alpha}\over
{T(\varphi, r)}}$

$\bullet$ for a sequence $(r_k)$ to be specified
$$T[\varphi]:= \lim_kT_{r_k}[\varphi].$$
\end{slide}

\begin{slide}
\par\vspace{\bigskipamount}\noindent
{\large\bf Basic properties of $T[\varphi]$}

\par\vspace{\bigskipamount}\noindent
$\bullet$ It is a closed, positive current.
\vskip 5pt \hskip 0.3cm $\alpha= d\beta$ \hskip 0.3cm $\Longrightarrow$ \hskip 0.3cm
$\int_1^r{{dt}\over {t}}\int_{\Delta(t)}\varphi^*\alpha= 
\int_1^r{{dt}\over {t}}\int_{S(t)}\varphi^*\beta$.
\par\vspace{\bigskipamount}\noindent
\hskip 0.3cm This last quantity is bounded by
$$C(\beta)\int_{\Delta(r)\setminus \Delta(1)}\vert \varphi^\prime (z)\vert _\omega
{{d\lambda}\over {\vert z \vert }}\leq C(\beta)\bigl( r\log r{{\partial}\over{\partial r}}T(\varphi, r)
\bigr)^{1/2}$$

\hskip 0.3cm We have \hskip 0.3cm $\displaystyle 
\lim\inf_{r\mapsto \infty} {{r\log r{{\partial}\over{\partial r}}T(\varphi, r)}\over {T(\varphi, r)^2}}= 0$

\par\vspace{\bigskipamount}\noindent
\noindent {\bf Examples}

$\bullet$ $z\mapsto (z, \exp(z))$, on ${\Pp}^1 \times{\Pp}^1$ 
\hskip 0.3cm $\Longrightarrow$  \hskip 0.3cm
$T[\varphi]= [\infty\times {\Pp}^1]$

$\bullet$ $z\mapsto (\exp(z), \exp(iz))$, on ${\Pp}^1\times{\Pp}^1$ 
$$T[\varphi]= [\infty\times {\Pp}^1]+ [0\times {\Pp}^1]+ [{\Pp}^1\times 0]+ [{\Pp}^1\times \infty]$$

\end{slide}

\begin{slide}
\par\vspace{\bigskipamount}\noindent
{\large\bf Numerical properties of $T[\varphi]$}
\par\vspace{\bigskipamount}\noindent

{\bf Proposition 1} {\sl Let $\{\alpha\}\in H^{1, 1}(X, \Rr)$ be a real cohomology class, 
which contains a positive 

\hskip 2.5cm current $T\in \{\alpha\}$. Then $(T[\varphi], \alpha)\geq 0.$}
\par\vspace{\bigskipamount}\noindent
{\bf Remark.} If $\{\alpha\}= c_1(D)$ where $D$ is an effective line bundle \hskip 0.3cm
$\Longleftarrow$ \hskip 0.3cm McQuillan.
\par\vspace{\bigskipamount}\noindent
{\bf Proof.} Based on

{\bf Theorem}{(Demailly)} {\sl Let $(X, \omega)$ be a compact complex manifold, 
and let $T= \alpha + i\partial\overline\partial f\geq 0$ be a 
closed positive current. Then for each $m\geq 1$, there exist
a function $f_m$, with logarithmic singularities, such that 
$T_m:= \alpha + i\partial\overline\partial f_m\geq - 1/m\omega.$}

As $\varphi$ is non-degenerate, the function $f_m\circ \varphi$ is not 
identically $-\infty$, and we have

$$\int_1^{r_k}{{dt}\over {t}}\int_{D(t)}\varphi^*(\alpha+ 1/m\omega )
+ \int_0^{2\pi} f_m(\varphi(r_ke^{i\theta}))d\theta/2\pi \geq {\cal O}(1)$$

\noindent We renormalise, and take the limit as $k\mapsto\infty$
and we get 
$$(T[\varphi], \alpha+ 1/m\omega)\geq 0$$ 

\end{slide}

\begin{slide}
\par\vspace{\bigskipamount}\noindent
{\large\bf Multiplicities of $T[\varphi]$}
\par\vspace{\bigskipamount}\noindent
Let ${\cal I}\subset {\cal O}_X$ be a coherent ideal sheaf

$\bullet$ $\exists$ quasi-psh function $\psi_{\cal I}$ on $X$ such that locally 
$$ \psi_{\cal I}(w)= 1/2 \log\sum_k\vert g_k(z)\vert ^2$$
\hskip 0.3cm (modulo smooth terms;  $(g_k)$ are the local generators of the ideal ${\cal I}$)

$\bullet$ we define
$${\rm mult} 
(T[\varphi], {\cal I})
:= \lim _{k\mapsto \infty}{{1}\over {T(\varphi, r_k)}}
\Bigl(\sum \nu_j\log {{r_k}\over 
{\vert z_j\vert}}+ 
\int_0^{2\pi}\bigl(\log\varepsilon-  \psi_{\cal I}( \varphi(r_ke^{i\theta}))\bigr)_+{{d\theta}\over {2\pi}}
\Bigr )$$

\hskip 0.3cm $(z_j)\subset \Cc$ are the zeros of the functions 
$(g_k(\varphi))$, with multiplicities $\delta_j$.

\par\vspace{\bigskipamount}\noindent 
$\bullet$ ${\cal I}= {\cal I}_{x_{0}}$ \hskip 0.3cm $\Longrightarrow$  \hskip 0.3cm
${\rm mult}(T[\varphi], {\cal I})= (T[\widehat \varphi], c_1(E))$ 

\hskip 0.3cm where $\widehat X\mapsto X$ blow-up at $x_0$, $E$ is the exceptional divisor.

\par\vspace{\bigskipamount}\noindent
$\bullet$ Existence of positive multiplicities for $T[\varphi]$ should 
imply better numerical properties.

\end{slide}

\begin{slide}
\par\vspace{\bigskipamount}\noindent
{\bf Back to the dimension 2} 
\par\vspace{\bigskipamount}\noindent
$\bullet$ Let $\varphi:{\CCP}\mapsto {S}$ be an entire curve on a surface $S$, 
${\rm mult}(T[\varphi], x_0)> 0$. 

\hskip 0.3cm If $\pi:\widehat S\mapsto S$
blow-up at $x_0$, then
$$\{T[\widehat \varphi]\}= \{\pi^*T[\varphi]- {\rm mult}(T[\varphi], x_0)E\}$$
\par\vspace{\bigskipamount}\noindent
\hskip 0.3cm Proposition 1 \hskip 0.2cm $\Longrightarrow$  \hskip 0.2cm
$\{T[\widehat \varphi]\}^2\geq 0$  \hskip 0.2cm $\Longrightarrow$  \hskip 0.2cm
$\{T[\varphi]\}^2\geq {\rm mult}(T[\varphi], x_0)^2$

\par\vspace{\bigskipamount}\noindent
$\bullet$ Consider the sequence $\pi_k:(S_{k+ 1}, x_{k+1})\mapsto (S_{k}, x_{k})$

\hskip 0.3cm where $S_{k+1}$ is the blow-up of $S_k$ at $x_k$, and $x_{k+1}$ belongs to the
exceptional divisor. 

\hskip 0.3cm Then
$$\sum_k{\rm mult}(T[\varphi_k], x_k)^2\leq T[\varphi]^2$$

\hskip 0.3cm In particular \hskip 0.2cm $\lim_k{\rm mult}(T[\varphi_k], x_k)= 0$
\end{slide}

\begin{slide}
\par\vspace{\bigskipamount}\noindent

\noindent {\bf Theorem} {\sl Let $X$ be a compact K\"ahler 
manifold, and let $\varphi:{\Cc}\mapsto X$ be a Zariski dense holomorphic curve.
If ${\rm mult}(T[\varphi], Y)> 0$ for some irreducible analytic set $Y\subset X$ of codimension equal to $p$, 
then for any nef class $\{\alpha\}\in H^{1, 1}(X, \Rr)$ such that $p+ \nu(\alpha)> n$
we have $(T[\varphi], \alpha)> 0$.}

\par\vspace{\bigskipamount}\noindent
\noindent {\bf Sketch of the proof} \hskip 0.3cm$Y= x_0$ 

$\bullet$ 
$T_\varepsilon= \alpha+ \varepsilon\omega+ i\partial\overline\partial f_\varepsilon$
\hskip 0.3cm $\Longleftarrow$  \hskip 0.3cm concentration of mass (theorem 1) 
$$\nu(T_\varepsilon, x_0)\geq C\varepsilon^{1-\nu(\alpha)/n}$$ 

$\bullet$ Let $\pi:\widehat X\mapsto X$ be the blow-up of $X$ at $x_0$
$$\pi^*T_\varepsilon\geq C\varepsilon^{1-\nu(\alpha)/n}[E]$$ 
\par\vspace{\bigskipamount}\noindent
$\bullet$ The curve 
$\widehat \varphi$ is still non-degenerate, so by the proposition 1, we get
$$(T[\widehat \varphi], \{\pi^*T_\varepsilon- C\varepsilon^{1-\nu(\alpha)/n}E\})\geq 0$$
\par\vspace{\bigskipamount}\noindent
$\bullet$
$(T[\varphi], \alpha+ \varepsilon\omega)\geq C\varepsilon^{1-\nu(\alpha)/n}{\rm mult}(T[\varphi], x_0)$ \hskip 0.3cm $\forall \varepsilon> 0$



\end{slide}

\begin{slide}
\par\vspace{\bigskipamount}\noindent
{\large\bf How about the Lelong numbers?}

\par\vspace{\bigskipamount}\noindent
$\bullet$ we have \hskip 0.2cm $\nu(T[\varphi], x_0)\geq {\rm mult}(T[\varphi], x_0)$ 

$\bullet$ even if \hskip 0.2cm $\nu(T[\varphi], x_0)> 0$ \hskip 0.2cm 
it may happen that \hskip 0.2cm $\{T[\varphi]\}^2= 0$ 
\par\vspace{\bigskipamount}\noindent
\noindent {\bf Theorem 2.} {\sl Let $X$ be a compact K\"ahler manifold, and let
$\varphi :{\Cc}{\mapsto X}$ be a Zariski dense entire curve. 

Assume that $ \hskip 0.5cm
\left\{
\begin{array}{ll}
\nu (T[\varphi], x_0)> 0 \hbox{ for some point } x_0\in X \\\
% \mbox{\hskip 0.5cm}\& \\
\chi_YT[\varphi]= 0, \hbox{ for any proper analytic set }Y\subset X
\end{array}\right. 
$

\noindent Then $(T[\varphi], \alpha)> 0$,
for any nef, non-trivial $(1,1)$--class $\{\alpha\}$.}
\par\vspace{\bigskipamount}\noindent
\noindent {\bf Remark} In general, we have $T[\varphi]= \sum\nu_j[C_j]+ R[\varphi]$
(Siu decomposition)
\par\vspace{\bigskipamount}\noindent
$\bullet$ Assume $\dim(X)= 2$ \hskip 0.2cm $\Longrightarrow$ \hskip 0.2cm
$T[\varphi]= R[\varphi]$ \hskip 0.2cm and \hskip 0.2cm $\nu (R[\varphi], x_0)> 0$

$\bullet$ $R[\varphi]\wedge R[\varphi]$ is well defined, and it dominates $\nu (R[\varphi], x_0)^2$

$\bullet$ In higher dimension, use -- ideas from comparison theorems for Lelong numbers,

\hskip 4.3cm -- the theorem 1


\end{slide}

\begin{slide}
\par\vspace{\bigskipamount}\noindent
{\large\bf Siu decomposition of $T[\varphi]$}

\par\vspace{\bigskipamount}\noindent
Consider $\varphi^\prime:{\Cc}\mapsto {\Pp}(T_X)$ the derivative of the curve $\varphi$. 

We have the 
following link between the currents $T[\varphi]$ and $T[\varphi^\prime]$

\par\vspace{\bigskipamount}\noindent
\noindent {\bf Proposition 2} {\sl If \hskip 0.2cm $T[\varphi]= \sum\nu_j[C_j]+ R[\varphi]$
\hskip 0.2cm then \hskip 0.2cm
$T[\varphi^\prime]= \sum\nu_j[\widehat C_j]+ \sum\tau_j[F_j]+ R[\varphi^\prime]$ 
\noindent where
$\widehat C_j$ is the canonical lifting of $C_j$, and $F_j$ are containd in the fibers of
$\pi: P(T_X)\mapsto X$.}

\par\vspace{\bigskipamount}\noindent
\noindent {\bf Corollary.} $\pi_*R[\varphi^\prime]=  R[\varphi]$.

\par\vspace{\bigskipamount}\noindent
\noindent {\bf Question.} Can one give a precise statement about the correspondance between
$$R[\varphi^\prime]\hbox { and } R[\varphi]?$$

\end{slide}


\begin{slide}
\par\vspace{\bigskipamount}\noindent
{\large\bf Main question}

\par\vspace{\bigskipamount}\noindent
{\bf Conjecture} {\sl Assume that $\chi_YT[\varphi]\not = 0$. 

\hskip 2.2cm Then there exist a non-constant 
entire curve $\varphi_Y:{\Cc}\mapsto Y$.}

\par\vspace{\bigskipamount}\noindent
\noindent {\bf Partial results}

$\bullet$ $\dim(X)= 2$, and $\varphi$ tangent to a holomorphic foliation
${\cal F}$ (Brunella, McQuillan)

$\bullet$ $T[\varphi]=  [C]$, where $C$ is an irreducible curve (J. Duval)


\end{slide}%

\end{document} 